\documentclass[12pt, a4paper]{article}

\begin{document}

\textbf{{\Large Peer review of Nathan Sever's Paper}}

\vspace{1cm}

Before the conceptual things, some small things that I noticed:

\begin{itemize}
	\item Page 2, paragraph 2, ``geographic''.
	\item Page 12, paragraph 2, maybe use human capital instead of intellectual capital? Not a huge deal.
	\item Page 13, paragraph 1, didn't quite understand the however there. Maybe there was something written before that was erased? Seemed a bit out of context.
\end{itemize}

First, I wanted to congratulate you on the progress made so far. I was truly impressed when I opened up the file and I saw the amount of pages and work that were in there. Really good!

The literature review and dialogue with other papers throughout the whole analysis is exceptional. For any further analysis, I would recommend you keep doing it like this.

I also appreciate the contextualization of your research question and the overall regulatory landscape surrounding your research question. As an international student, it was easy to follow and interesting.

Regarding the models, I find that the idea of creating this synthetic control state for the conterfactual Minessota to be a great idea. I didn't know that this was done in diff-in-diffs regressions, but the concept is really well explained in the paper itself. For the seasonality issue, I found that part a bit hard to understand. From what I understand, you will be using the residual created by subtracting the Predictor variable from the true variables. Are you only going to do that with earnings? I think it will kill some of the heterogeneity and decrease sample size that you have in your data and, given that only few years worth of data are available after the treatment, it might dim even more the results (but I might be wrong). Also, how does that work for the model you proposed? Isn't the outcome regression related to a specific quarter? Doing a bit of research, I found that your approach is reasonable, but there are some more things to consider. My suggestion is a different approach (I might be wrong). I suggest that you split your year fixed effect into 4 different fixed effects, which will represent each quarter. This might help, but I am not sure.

For the interpretation part, I have some concerns about the start up interpretaion on page 12. Even though your argument makes some sense (smaller firms with NCAs are likely to be startups), it would be good to see some actual data that confirms this. It is not obviously true at first sight. And even if data confirms it, I think you would still be commiting some kind of ecological fallacy (see Robinson (1950), found this paper during this class :-) ). Other than that, I think they make sense.

Talk with me if you have any questions! Inspiring to see a peer that has done so much in the class. Would be happy to talk to you and brainstorm some ideas.



\end{document}